% !PRINTED 285
\weberpage{285}

%<<CONT
so ist nach \S.\,83, V. die Anzahl der Zeichenwechsel in
\weq{2}{1,\ \varDelta_1(\alpha), \varDelta_2(\alpha) \dots \varDelta_n(\alpha)}
gleich der Anzahl $N_\alpha$ der negativen Quadrate in $H$, und wenn wir also die entsprechende Zahl $N_\beta$ in
\[
1,\ \varDelta_1(\beta), \varDelta_2(\beta) \dots \varDelta_n(\beta)
\]
abzählen, so ist $N_\beta - N_\alpha$ die Anzahl der zwischen $\alpha$ und $\beta$ gelegenen Wurzeln von $f(x)$.

Für die cubische Gleichung ergeben sich die Ausdrücke aus dem Schluss des vorigen Paragraphen. Wir bilden die Kette
\[
a_0,\quad \frac{1}{a_0} H_{2,2},\quad \frac{1}{a_0}\left(H_{2,2} H_{1,1} - H_{1,2}^2\right),\quad \frac{1}{a_0} \varDelta,
\]
die offenbar dieselbe Anzahl von Zeichenwechseln hat wie (2), und erhalten so die vier Functionen
\[
\begin{aligned}
&a_0,\qquad -3 a_0 \alpha - a_1,\\
&2\alpha^2 a_0(a_1^2 - 3a_0 a_2) + \alpha(2a_1^3 + 9a_0^2a_3 - 7a_0a_1a_2)\\
&\quad + 3a_0a_1a_3 + a_1^2a_2 - 4a_0a_2^2,\qquad -f(\alpha)\,D.
\end{aligned}
\]
Nehmen wir zur Probe $\alpha = -\infty$, $\beta = +\infty$, so erhalten wir, wenn noch $a_0 = 1$ gesetzt wird, die Zeichenbestimmung
\[
\begin{aligned}
1, &\quad +1, \quad (a_1^2 - 3a_2), \quad D,\\
1, &\quad -1, \quad (a_1^2 - 3a_2), \quad -D.
\end{aligned}
\]
Man muss bei der Abzählung beachten, dass, wenn $D$ positiv ist, $a_1^2 - 3a_2$ nicht negativ sein kann [\S.\,47, (8)].

\begin{center}
\S.\,93.

\gesperrt{Grundzüge der Charakteristikentheorie}\footnote{Monatsberichte der Berliner Akademie, März und August 1869, Februar 1873, Februar 1878.}).
\end{center}

Die \gesperrt{Sturm}'schen Sätze werden in ausserordentlicher Weise durch die Charakteristikentheorie von \gesperrt{Kronecker} verallgemeinert, die dasselbe Ziel wie der \gesperrt{Sturm}'sche Satz für Gleichungssysteme mit beliebig vielen Veränderlichen verfolgt. Wir beschränken uns hier auf die Betrachtung des einfachsten Falles, den wir zur Einschliessung der complexen Wurzeln einer Gleichung anwenden wollen; auch bedienen wir uns hier unein-
%CONT>>
